% Level 2, Week 10 -- Project Studio 1
\lesson{10}{Project Studio 1}{Plan first, then build the core: a small program that works beats a big one that does not.}
\lessoninfo{Concepts: pseudocode and flowcharts (I), decomposition (R), core loop (R), all Level 2 (R/E)}{Tools: Thonny, paper, sticky notes, the Project Brief}{CSTA: 2-AP-10, 2-AP-13, 2-AP-18, 2-IC-22}

\begin{objectives}
  \item The student can turn a project idea into pseudocode or a flowchart and point to the loop, the decisions and the data in it. (L2.G8)
  \item The student can break the project into named functions and split the tasks fairly with a partner on a task board. (L2.G3, L2.G8)
  \item The student can build a core loop that runs without crashing: a game timer loop or an art-studio menu loop (Milestone M2). (L2.G2, L2.G7)
  \item The student can give kind, specific code-review feedback and act on feedback from a partner. (L2.G8)
\end{objectives}

\begin{materials}
  \item Print: workbook chapter 10 and homework 10 (one per student). Each team brings its Project Brief and its Milestone M1 sheet from week 9.
  \item Paper, pencils and a few coins (paper prototype); per team a task board: a sheet folded into three columns labelled \textbf{TO DO}, \textbf{DOING}, \textbf{DONE}, plus 10 sticky notes.
  \item On each Mac: Thonny open and a folder \texttt{Documents/my-project} per team (Finder, \keys{Cmd}+\keys{Shift}+\keys{N}). Teams may copy the files from \texttt{2-intermediate/project/starter/} into it, or start from a blank file.
  \item A timer for 8-minute swaps. Prep: read this plan (10 minutes). Milestones: M1 plan (week 9) → \textbf{M2 core loop (today)} → M3 data and tests, M4 polish (week 11) → Showcase (week 12); full rubric in the Project Teacher Guide.
\end{materials}

\needspace{16\baselineskip}
\section{At a glance}
\begin{timeplan}
  \timeslot{0--8}{Warm-up}{Paper prototype: each team acts out one minute of its game or studio. Check Milestone M1 (pseudocode or flowchart + task split).}
  \timeslot{8--18}{Concept}{From plan to code: name the functions, build the \enquote{skateboard} first, split tasks fairly on a task board.}
  \timeslot{18--28}{Guided demo}{You type a core-loop skeleton from pseudocode: settings dictionary, state dictionary, one logic function, keys, timer loop.}
  \timeslot{28--52}{Hands-on build}{Teams build their own core loop (M2). Swap every 8 minutes. Minute 44: partner code review. Last 4 minutes: free-build personal touch.}
  \timeslot{52--60}{Wrap-up}{Each team shows its core running (30 seconds). Tick M2. Exit ticket, reflection, homework 10.}
\end{timeplan}


\section{Warm-up (0--8 min)}
\begin{unplugged}[Paper prototype]
Each team draws its screen on paper (player, falling things, score; or the art-studio menu) and
uses coins for the moving parts. One student is \textbf{the computer}: each time you clap is one
\enquote{tick} of the timer, so the computer moves every falling coin down one finger-width and
checks for touches. The other student is \textbf{the player} and calls \enquote{Left!} or
\enquote{Right!}. Art Studio teams hold up menu cards (\textbf{Draw}, \textbf{Save}, \textbf{Load},
\textbf{Quit}) and the computer does exactly what the card says.
After one minute ask: \enquote{What did the computer have to remember?} Those answers are the
variables, the list and the dictionary. With 1 student, you are the computer.

\textbf{M1 check} (1 minute per team): the plan shows a loop, at least one decision
(\enquote{if the star touches the basket}), the data (a list, a dictionary, a file) and at least
three function names. Add anything missing in pencil now.
\end{unplugged}

\section{Concept (8--18 min)}
\begin{teacherscript}
“Today your plan becomes a program that runs. Game makers build the smallest version that works
first: a window, a player that moves, a loop that keeps going. No score, no fancy colours yet. Then
every new feature is small and safe. Look at your pseudocode. Each line that does one job becomes a function. \enquote{Move the player
left but not off the screen} is a function. What goes in? The x position and the direction. What
comes out? The new x. Parameters in, return value out, and next week we can test it. Now the team part. Write every function name on its own sticky note and put them in TO DO. Take
turns choosing, one note each, like picking teams fairly. Nobody gets only the boring jobs, nobody
gets only the fun ones. Started a task? Move it to DOING. Works? Move it to DONE.”
\end{teacherscript}

\begin{analogy}[A skateboard before a car]
If you want to build a car, do not build one wheel, then another wheel, then the doors: nobody can
ride a wheel! Build a skateboard first. It is tiny, but you can ride it. Then a scooter, a bike,
the car. Every version works. Our core loop is the skateboard.
\end{analogy}

Then open workbook chapter 10 at the flowchart: ovals are start and end, rectangles are steps,
diamonds are yes/no questions. Trace one lap of the loop with your finger.

\section{Guided demo (18--28 min)}
\begin{workedexample}[From pseudocode to a walking skeleton]
In Thonny: \menu{File > New}, then \keys{Cmd}+\keys{S} and save as \texttt{mygame.py} in
\texttt{Documents/my-project}. Type the program below in this order, pressing \keys{F5} after each
step so the students see that every step works:
\begin{enumerate}
  \item The pseudocode as \texttt{\#} comments first: \enquote{set up}, \enquote{every tick: move, check, score}.
  \item \texttt{SETTINGS} (numbers we may change later) and \texttt{state} (things that change while playing): two dictionaries.
  \item \texttt{clamp} and \texttt{next\_x}: pure logic, no turtle. Say: \enquote{This is the brain. Next week we test the brain.}
  \item The key handlers and \texttt{game\_loop}. Point at this line:\newline \texttt{screen.ontimer(game\_loop, 50)}\newline \enquote{Run me again in 50 milliseconds. No brackets after \texttt{game\_loop}: we hand the function over, we do not call it now.}
  \item The start-up code under \texttt{if \_\_name\_\_ == "\_\_main\_\_":} so tests can import this file without opening a window.
\end{enumerate}
Press \keys{F5}, click inside the turtle window, press \keys{←} and \keys{→}: the square slides and
stops at the edges. That is M2 for an arcade game. Art Studio teams build a \texttt{while} menu loop
instead (tested example in the answer key, week 10; the turtle window looks frozen while Python waits
for \texttt{input()}: answer in Thonny's Shell).
\end{workedexample}

\begin{pycode}[mygame.py -- core loop skeleton (tested)]
# My game - core loop skeleton (Milestone M2)
import turtle

SETTINGS = {"step": 20, "left_edge": -280, "right_edge": 280}
state = {"score": 0, "running": True}

def clamp(value, low, high):
    """Keep value between low and high."""
    return max(low, min(high, value))

def next_x(x, direction):
    """Where the player moves to: direction is -1 (left) or 1 (right)."""
    new_x = x + direction * SETTINGS["step"]
    return clamp(new_x, SETTINGS["left_edge"], SETTINGS["right_edge"])

def go_left():
    player.setx(next_x(player.xcor(), -1))

def go_right():
    player.setx(next_x(player.xcor(), 1))

def game_loop():
    if not state["running"]:
        return
    # TODO: move falling things, check for catches, add to the score
    screen.update()
    screen.ontimer(game_loop, 50)   # run again in 50 milliseconds

if __name__ == "__main__":
    screen = turtle.Screen()
    screen.tracer(0)
    player = turtle.Turtle("square")
    player.penup()
    player.goto(0, -250)
    screen.listen()
    screen.onkey(go_left, "Left")
    screen.onkey(go_right, "Right")
    game_loop()
    turtle.done()
\end{pycode}

\section{Hands-on build (28--52 min)}
\begin{pairroles}[Driver, navigator and the task board]
The \textbf{driver} types. The \textbf{navigator} holds the plan, reads the next step aloud, spots
typos and moves the sticky notes. Swap when the timer rings and press \keys{Cmd}+\keys{S} before
every swap. \textbf{1 student:} you are the navigator; ask \enquote{what does the plan say next?}
\textbf{2:} one team. \textbf{3:} a trio with driver, navigator and \textbf{tester} (runs the program
after each change), rotating all three roles. \textbf{4:} two teams; you float between them.
\end{pairroles}

Goal by minute 44, \textbf{M2}: arcade teams have a player that moves with keys and one thing that
moves by itself on the timer; Art Studio teams have a menu loop that draws something and quits
cleanly. Teams may remix the week 9 \enquote{Dodge the Meteors} exercise
(\texttt{excercises/week09/ex-dodge-meteors/}) or pieces of Star Catcher, with a credit comment at
the top, for example \texttt{\# Movement code remixed from Dodge the Meteors}.

\begin{bughunt}[Partner code review (minute 44, 6 minutes)]
Teams swap seats (with one team, you and the student swap). The reviewer runs the other team's code
and reads it with this checklist:
\begin{checklist}
  \item Does it run? If there is a red error, which line does the traceback point to?
  \item Are names clear? (\texttt{next\_x} good, \texttt{thing2} not so good.)
  \item Does each function do one job? Are numbers you might change in a settings dictionary?
\end{checklist}
Look out for the three classic week 9 bugs: \texttt{screen.onkey(go\_left(), "Left")} (the brackets
run the function once now, not on the key press), a missing \texttt{screen.listen()}, and a
\texttt{game\_loop} that never calls \texttt{ontimer} again (things move once, then freeze).
The reviewer writes \textbf{two stars and a wish} on a sticky note. The authors decide what to do
with the wish: it is their project.
\end{bughunt}

\begin{freebuild}[Personal touch (last 4 minutes)]
Add one thing that makes it yours: a window title (\texttt{screen.title("Comet Catch")}), a palette
background colour, or a player shape (\texttt{"turtle"}, \texttt{"circle"}, \texttt{"triangle"}).
\end{freebuild}

\section{Wrap-up (52--60 min)}
Each team runs its program for 30 seconds; clap for every working loop, however small. Tick M2 on
the Project Brief (or write \enquote{M2: first 10 minutes next week}). Hand out homework 10.

\begin{formativecheck}[Exit ticket]
\textbf{Look for:} the program runs; something repeats on its own (moves on a timer, or the menu comes
back); at least two functions written by the team; sticky notes in DONE.
\textbf{Ask:} “In \texttt{screen.ontimer(game\_loop, 50)}, what does 50 mean, and why is there no
\texttt{()} after \texttt{game\_loop}?”
\textbf{Answer:} 50 milliseconds until the loop runs again. We hand the function to turtle so it can
call it later; with brackets it would run once, straight away.
\end{formativecheck}

\begin{reflection}
What was the smallest version of your project that worked today? Which sticky note are you proudest
of? Was the task split fair, and what would you change next week?
\end{reflection}

\begin{safety}[Collaboration: splitting tasks fairly]
Fair teams share the interesting and the boring jobs, swap the keyboard on time and give credit.
Each team puts both nicknames in a comment on line 1 (\texttt{\# Made by Comet and Pixel}) and a
credit comment for any code it remixed. Say: \enquote{Using someone's idea is fine. Hiding that you
used it is not.} Real names are not needed anywhere in the code.
\end{safety}

\section{Notes for the teacher}
\begin{commonmistakes}
  \item \textbf{Fancy parts first} (title screen, colours) before anything moves. Point at the skateboard: \enquote{What is the smallest thing that can run?}
  \item \textbf{A file named \texttt{turtle.py} or \texttt{random.py}.} Python imports the team's file instead of the real module and odd errors appear. Rename it (\texttt{mygame.py}).
  \item \textbf{Keys do nothing.} Click inside the turtle window, check \texttt{screen.listen()}, and check key names: \texttt{"Left"} with a capital L, \texttt{"space"} in lower case.
  \item \textbf{One student hogs the keyboard.} Keep the timer strict and say \enquote{swap!} cheerfully.
\end{commonmistakes}

\begin{differentiation}
\supportpath Start from \enquote{Dodge the Meteors} or the skeleton above and change one thing at a
time (shape, speed, what falls). Shrink the plan: one falling object, two keys, score later. A remix
with credit is a valid project start.
\extensionpath Keep several falling things in a list and move them with one \texttt{for} loop, or
(Art Studio) choose between shape types with a dictionary of drawing functions.
\end{differentiation}

\begin{stretchchallenge}[Game over and play again]
When \texttt{state["running"]} becomes \texttt{False}, write \enquote{GAME OVER -- press R} on the
screen. Pressing \keys{r} resets the score, sets \texttt{running} back to \texttt{True} and calls
\texttt{game\_loop()} again.
\end{stretchchallenge}

\begin{teachernote}[If you are short on time]
Cut the paper prototype to 3 minutes. Never skip the M2 run-through at the end: it is the most
motivating moment of the day. The code review can move to the start of week 11.
\end{teachernote}
